\subsection{Experimental Design for Ecological Evaluation}\label{sec:protocol}
The following protocol is proposed for the observations to be added later.
It is not a description of an already completed ecological benchmark.

\subsubsection{Research questions and controls}
The primary questions are whether trained policies improve hunting/feeding and
survival on unseen worlds; whether paired inheritance preserves useful behavior;
and how sensing, capacity, terrain and reproduction affect those outcomes.
Each study should specify its primary endpoint before inspecting results.

Separate training, validation/selection and held-out test seeds. Freeze selected
policies and evaluate them against fresh untrained and random-action controls
under matched physical configurations. For asymmetric comparisons, freeze the
opponent to a documented baseline or opponent pool. Otherwise improvement in
one species can reflect deterioration in the other. Training-time learning curves
and frozen-policy generalization should be reported separately.
Frozen evaluation disables birth mutation as well as learning, while parental
crossover still initializes newborns. Mutation ablations therefore require
explicit training conditions and a clearly stated evaluation birth policy.

Model selection must use training/validation information only. Record source
experiment, numbered archive, species/sex criterion, signature and selected
artifact identity. A high score obtained through repeated seed search is not an
independent test result. Best-so-far retention and early stopping are selection
mechanisms whose settings belong in the experimental description.

\subsubsection{Comparisons and budgets}
Useful focused ablations include $k=1$ versus $k=4$; matched versus asymmetric
predator capacity/sight; mating features absent/present; parental crossover
with mutation disabled/enabled; independent initialization versus chained
champions; and lifespan versus balanced/lineage selectors. Keep nonstudied body,
resource and reward coefficients fixed. Incompatible observation/action/network
contracts require fresh training rather than forced weight loading.

For geography studies, include flat, water and elevation/enclosure cases with
reported map and spawn layout. Report population caps, initial energy grants,
health, food growth, costs, reproductive windows and episode horizon. A population
at its cap is censored by that rule and cannot be interpreted as unconstrained
reproductive success.

Report both real-world animal-steps and wall-clock cost. For GRPO, also report
trial groups, group size, horizon, scheduling and hypothetical simulation steps.
Equal episode counts can hide very different population sizes and update work.
Compare learning budgets separately from native implementation throughput.

\subsubsection{Metrics and statistical unit}
Let $A_z=\sum_t N_z(t)$ denote animal-steps for species $z$. Rates should state
their denominator, for example kills per 1,000 predator animal-steps or successful
feeds per attempted feed. Record extinction frequency/time, population curves,
births, predation/starvation/drowning deaths, completed lifespans, survivor ages,
energy intake/demand and viable generations. Separate newborn grants from food
when calculating energy coverage.

An independent seeded run is the primary replication unit. Animals and frames
within a run interact and are not independent samples. Report per-seed outcomes,
number of seeds, effect sizes and uncertainty using a stated estimator, such as
a paired bootstrap over seed-level differences. Predefine the seed count/budget
rather than stopping after an attractive result. Treat living final agents as
right-censored; use an appropriate survival analysis or label descriptive
age-at-death summaries without implying unbiased lifetime estimates.

Diagnostics are sampled. Q values, action probabilities and entropy describe
model outputs on changing observations, not prediction accuracy. Missing action
attempts have undefined success rates, not a zero-success observation. Analyses
requiring consecutive frames must disclose coverage, terminal-action gaps and
recording cadence. Report missing or unidentifiable quantities explicitly.
